% Owner: Kashaf Ali / Member A. Complete rewrite authorized 5 October 2026.
% Local font encoding restores the template's Times regular/bold faces under XeTeX.
\begingroup
\renewcommand{\encodingdefault}{T1}
\fontencoding{T1}\selectfont
\chapter{Introduction}
A patient's clinical history develops across encounters, as consultations add information relevant to later assessments. Symptoms described during a visit, laboratory observations, medical images and medication information represent different parts of that history. Considering these inputs together requires attention to their source, timing and relationship to the current case. A record from an earlier encounter may provide useful context, but its meaning cannot be separated from when it was recorded or why it was collected. Electronic health record research illustrates the importance of organising clinical data through identifiable patient and encounter relationships.\cite{zeestA:mimicIV} Zeest addresses this information-management problem through longitudinal patient memory, linking the current consultation with relevant historical information. Its focus is the continuity of the patient case across encounters rather than the interpretation of an isolated input. This perspective establishes the need for a workflow in which the physician can inspect both present information and the context in which earlier information was obtained.

Consultation documentation adds another task to reviewing a clinical case. The spoken exchange between a physician and patient contains information that must be represented clearly before it can be used alongside other records. In Zeest, real-time transcription is intended to convert consultation audio into structured clinical notes, giving the current encounter a documented representation that can enter the wider workflow. The physician needs to review these notes, correct transcription errors and understand which content came from the consultation. Imaging, laboratory information and medication context also require suitable representations, since they cannot be treated as interchangeable text inputs. Connecting them to the same patient and encounter supports subsequent analysis. The central requirement is therefore more than collecting data in one place: the system must preserve the relationships that allow a physician to interpret the information coherently. Consultation capture, note review and the handling of additional inputs form a connected physician workflow.

Language models offer a way to work with clinical text, but their capabilities must be interpreted against the tasks on which they are evaluated. Nori and colleagues examined GPT-4 using medical examination materials and benchmark datasets, while also discussing accuracy and safety challenges.\cite{zeestA:noriMedical} Such work motivates investigation of language-based assistance without establishing that an examination result demonstrates suitability for an individual patient's care. Retrieval-augmented generation provides a complementary approach by combining generation with access to external passages, making retrieved material available to the reasoning process.\cite{zeestA:rag} For Zeest, the design question is how these capabilities can be connected with patient-specific memory and specialised clinical processing. A generated recommendation must be distinguishable from its supporting sources, and the physician must inspect both. The integration must also preserve the role of imaging and laboratory findings instead of relying on text alone. These requirements motivate a coordinated workflow with evidence retrieval, separate review and explicit physician judgement.

Zeest is an agentic medical copilot designed to bring longitudinal memory, consultation documentation and multimodal clinical analysis into that workflow. The Patient Memory Agent retrieves relevant history, while the Clinical Case Agent organises the current inputs and historical context. Imaging, Laboratory and Medication Safety agents contribute findings within their responsibilities. Existing validated diagnostic models and established knowledge resources provide the diagnostic groundwork; the project's contribution lies in how these capabilities are integrated with memory, reasoning and review. The Evidence and RAG Agent supplies relevant guidelines with citations, and the Clinical Reasoning Agent produces a draft diagnosis and treatment plan. Before presentation, the Safety and Critic Agent reviews the generated content. The physician then approves, modifies or rejects the draft, and only explicitly approved decisions enter the clinical record. An audit history records the review actions. This arrangement supports investigation of the contributions of memory, agent organisation, critic review and multimodal inputs while keeping the clinical decision under physician control.

\section{Purpose of this Document}
Zeest aims to develop a clinical decision-support system that maintains patient context across encounters and connects it with the information collected during a consultation. The objective is to support a reviewable assessment by combining transcription, specialised agents, existing diagnostic models and retrieved evidence. The research examines four related questions: whether longitudinal memory improves recommendation quality, whether coordinated agents perform better than a single language model, whether critic review reduces unsupported outputs, and whether multimodal inputs improve recommendations compared with text-only information. These questions concern the contribution of the integrated workflow rather than the development of a new diagnostic model. They require comparisons that preserve the relationship between patient case, the available inputs and the reference used to assess the output. The purpose of the document is to establish an account of the problem, system requirements and design decisions needed to investigate these questions. It provides the basis for developing the system and defining a controlled evaluation of its components.

The investigation connects the application's clinical workflow with a comparative research framework. Patient and encounter information must remain identifiable across the processing stages so that differences between system variants can be assessed on the same cases. A single-language-model baseline, a memory-augmented variant and the full Zeest system form the central comparison, with critic and multimodal contributions examined through appropriately controlled contrasts. The requirements identify the operations each variant must support, while the methodology and design establish how those operations are connected. The specification defines the physician's role in reviewing evidence and controlling the final recorded decision. Clinical domain coverage depends on suitable data and existing validated models, and development uses de-identified or synthetic information. These boundaries keep the investigation focused on integration and decision support. The current deliverable establishes the requirements and design foundation; implementation and evaluation will provide the evidence needed to answer the research questions rather than assuming a favourable outcome in advance.

\textcolor{red}{PLANNED: System implementation and comparative evaluation will follow the agreed requirements and methodology; no experimental results are presented.}

\section{Intended Audience}
Academic evaluators and the project supervisor need to understand Zeest, the research contribution and its objectives and design. The student development team needs a consistent specification of user actions, processing responsibilities and information that must be retained across encounters. Researchers and software developers interested in clinical language models, retrieval and agent integration can use the same account to examine how the components cooperate and where the research questions arise. Physicians are an audience because the system is designed around their review of patient context, clinical findings and generated recommendations. For these readers, the explanation of evidence presentation, critic findings and approval controls must be understandable without knowledge of the underlying orchestration frameworks. The writing therefore connects technical concepts with the clinical workflow they support. This combination of academic, technical and professional readers guides the detail: requirements must be precise enough for development while remaining clear about the physician's interaction with the system.

\section{Definitions, Acronyms, and Abbreviations}
Patient history, current encounter information and retrieved evidence serve distinct roles in Zeest. Longitudinal memory refers to patient-specific information retained across encounters, whereas evidence retrieval concerns external clinical material used to support the reasoning process. An encounter identifies the context in which consultation notes and other inputs are collected. A clinical draft is the generated diagnosis and treatment plan awaiting physician review, and an approved decision is the content explicitly accepted for entry into the clinical record. An audit entry records a review action; its presence does not by itself mean that the associated draft was approved. These distinctions are particularly important when describing storage, retrieval and modification, since each operation affects a different part of the workflow. Table~\ref{tab:kashaf:terms} defines the principal abbreviations and terms used in the following chapters. The terminology separates the information supplied to the system, the findings generated by its components and the decision made by the physician.

\begin{table}[htbp]
\centering
\small
\caption{Abbreviations and terms used to describe Zeest and its clinical workflow.}
\label{tab:kashaf:terms}
\begin{tabular}{|p{0.24\textwidth}|p{0.66\textwidth}|}
\hline
Term & Meaning \\ \hline
AI & Artificial Intelligence. \\ \hline
LLM & Large Language Model. \\ \hline
CDSS & Clinical Decision Support System. \\ \hline
RAG & Retrieval-Augmented Generation. \\ \hline
SDG & Sustainable Development Goal. \\ \hline
Encounter & An identifiable consultation or visit associated with a patient. \\ \hline
Longitudinal memory & Patient-specific information retained and retrieved across encounters. \\ \hline
Hybrid memory & Relational records combined with semantic vector representations. \\ \hline
Multimodal input & Clinical information represented through notes, images, laboratory information and medication context. \\ \hline
Agent & A specialised processing component with a defined responsibility. \\ \hline
Clinical draft & A generated diagnosis and treatment plan awaiting physician review. \\ \hline
Approved decision & Clinical content explicitly accepted by the physician for recording. \\ \hline
Audit history & Recorded physician actions associated with review and decision making. \\ \hline
Ablation study & A comparison of system variants used to examine component contributions. \\ \hline
\end{tabular}
\end{table}

\section{Conclusion}
Zeest brings together the processing needed to review a patient case with continuity across encounters. Its central contribution is the integration of consultation transcription, hybrid patient memory, specialised analysis and retrieved evidence within a physician-controlled workflow. The generated diagnosis and treatment plan remains a draft until the physician makes an explicit decision, and the audit history preserves the relationship between that draft and the review actions. This distinction shapes both the user-facing requirements and the underlying data design. The research questions examine how memory, agent coordination, critic review and multimodal information contribute to its recommendations. Addressing those questions requires a specification before development and a controlled comparison after implementation. The introduction therefore establishes the clinical context, the investigation and the terminology needed to follow the remaining technical discussion. The next chapters develop these ideas into a defined scope, an assessment of related work and a detailed description of system behaviour and design.

Chapter~1 introduces the clinical information problem, the purpose of Zeest and the terms used to distinguish patient context, generated drafts and physician decisions. Chapter~2 develops the project vision, including the problem statement, objectives, scope, constraints, SDG alignment and stakeholder needs. Chapter~3 examines relevant research and related applications to explain the technical background and the limitations that motivate the investigation. Chapter~4 specifies the features, functional and non-functional requirements, physician use cases, resource needs, interface coverage, database relationships and project risks. Chapter~5 describes the proposed approach and methodology, connecting the agent workflow with the controlled comparisons used to investigate the research questions. Chapter~6 presents the high-level and low-level design, including architecture, subsystem responsibilities, domain concepts and review policies. The References section identifies the sources supporting the technical and research discussion, while the Appendix provides supporting material that complements the main chapters. Together, these parts connect the clinical motivation to the requirements and design needed for development and evaluation.

\endgroup
